\section{Conclusion}
\label{sec:conclusion}

We presented an adaptive spectral defense framework for protecting
deep-learning-based automatic modulation classifiers against adversarial attacks.
Our framework requires no retraining of the classifier and incurs sub-millisecond
overhead, making it suitable for always-on deployment in real-time spectrum
monitoring and cognitive radio applications.

Through systematic evaluation on the RML2016.10a dataset against five attacks
(CW, EAD-L1, EAD-EN, FGSM, PGD) and comparison with five calibrated classical
signal processing baselines, we established four principal findings:

\begin{enumerate}

\item \textbf{Adversarial attacks on AMC are control-plane attacks.}  CRC
  integrity analysis shows that adversarial perturbations collapse AMC accuracy
  to near zero while leaving the underlying signal data largely intact (CRC pass
  rate $> 71\%$ with oracle demodulation).  This characterization motivates
  recovery as the correct defense strategy: the perturbation need only be
  removed to restore correct classification.

\item \textbf{Adaptive-K FFT filtering outperforms all classical baselines on
  optimization-based attacks.}  By automatically selecting the number of
  retained spectral components per sample based on cumulative energy threshold
  ($\eta = 0.95$), Adaptive-K achieves 77.4\% weighted-average accuracy against
  CW (vs.\ 72.7\% for the best classical filter, Kalman) and 79.8\% against
  EAD-L1 (vs.\ 72.5\% for Kalman)---improvements of 4.7 and 7.3 percentage
  points, respectively.

\item \textbf{Classical signal processing filters are ineffective against
  optimization-based attacks despite per-SNR calibration.}  Savitzky-Golay,
  Wiener, and Kalman filters all perform at or below the undefended baseline on
  CW and EAD attacks, because their smoothing operations spread spectrally
  concentrated adversarial energy rather than removing it.  This negative result
  is practically important: it demonstrates that deploying conventional filtering
  as a security measure provides a false sense of protection against
  optimization-based adversaries.

\item \textbf{Spectral-gated routing improves recovery for wideband analog
  modulations.}  By bypassing Top-K filtering for samples whose spectral
  flatness exceeds a threshold (indicating wideband energy distribution),
  Spectral-Gated avoids distorting AM-SSB and WBFM signals where Top-K alone
  introduces classification-relevant artifacts.  Spectral-Gated achieves
  76.1\% on CW vs.\ 77.4\% for Adaptive-K, offering a useful accuracy--overhead
  trade-off for heterogeneous signal environments.

\end{enumerate}

Several directions remain open for future work.  First, the current evaluation
assumes a defense-unaware attacker; an adversary with knowledge of the Top-K
operator could craft perturbations that survive the frequency mask.  Evaluating
such \emph{adaptive attacks}---where the attack optimization incorporates the
defense---is an important robustness validation that we leave to future work.
Second, we focused on RML2016.10a (128-sample signals, 11 classes); extending
the evaluation to RML2018.01a (1024-sample signals, 24 classes) would establish
the generality of the spectral defense approach across longer signals and a
richer modulation set.  Third, combining adaptive spectral filtering with
adversarial training could provide complementary robustness: training-time
regularization hardens the classifier's decision boundaries while
inference-time filtering removes perturbation energy before the signal reaches
the network.
